\chapter{Conclusion}
\label{ch:conclusion}

The world-model calculus is not a replacement for PLN\@.  It is the
mathematical foundation that PLN always needed.

The central insight is that PLN's original instincts were largely
correct: evidence counts as the fundamental representation, revision
as the primary operation, and link-level rules as the practical
inference mechanism.  What was missing was the explicit separation of
layers and the honest statement of assumptions.  The passage from
intuition to formalization required not new ideas so much as the
discipline to say precisely what each idea means, what it assumes, and
where it breaks.

\section{What the Refoundation Gives Us}

The calculus provides five things that the original PLN lacked:

\begin{enumerate}
\item \textbf{One kernel.}  The world-model posterior-state calculus,
  with evidence compositionality as the fundamental law
  (Chapter~\ref{ch:world-models}).  Every inference rule, every view
  projection, every scaling result derives from this single interface.

\item \textbf{Many views.}  Strength, confidence, intervals,
  distributions---all as typed projections from the kernel
  (Chapter~\ref{ch:views}).  The views are not competing
  representations; they are different questions asked of the same
  underlying state.

\item \textbf{Explicit side conditions.}  Every compiled rule carries
  its assumptions (Chapter~\ref{ch:sigma-rewrites}), with
  counterexamples showing what happens when they are violated.  The
  d-separation typeclass (Chapter~\ref{ch:bn-models}) automates the
  checking, but the assumptions are always visible.

\item \textbf{Cleaner boundaries.}  The no-go theorem
  (Chapter~\ref{ch:no-go}), the identifiability failures, the pooling
  non-uniqueness, the coverage surrogate limitations, the five
  certified-chaining impossibilities
  (Chapter~\ref{ch:scale-control})---all stated precisely and proved
  formally.  These are not embarrassments; they are the most valuable
  theorems in the development, because they tell you where not to
  waste effort.

\item \textbf{Beyond subsumption.}  The bridges to neighboring
  frameworks are not merely compatibility results.  The
  compatibility-gated experiment demonstrates that decomposing an
  MLN's clause-weight signal into interpretable world-model evidence
  roles yields substantial improvement---not by adding information,
  but by factoring the information that was already there.
\end{enumerate}

\section{The Shape of the Result}

The result is a PLN that knows what it knows, states what it assumes,
and marks what it cannot do.  That is not a weakness.  That is the
beginning of a trustworthy reasoner.

The formalization is not complete---Chapter~\ref{ch:limitations}
states what remains undone, and Chapter~\ref{ch:future} sketches the
natural continuations.  But the kernel is stable, the evidence algebra
is settled, the compiled inference rules are proved exact under their
stated conditions, and the scaling theory identifies both the
achievable guarantees and the fundamental impossibilities.

The deepest insight is simple: inference systems should not confuse the
map with the territory.  A truth value is a projected summary of
evidence, not the evidence itself.  A probability is a query result,
not what the system carries.  What the system carries is a revisable
state of knowledge.  What the system produces is an extraction from
that state.  What practitioners see is a view projection of that
extraction.

This distinction has practical consequences.  When evidence is primary:
\begin{itemize}
\item Revision becomes compositional (new evidence adds to old
  evidence)
\item Queries become principled (extraction is a homomorphism)
\item Views become interchangeable (the same evidence supports
  multiple operational summaries)
\item Inference becomes modular (local computations compose globally)
\item Uncertainty becomes transparent (we track what we know, not
  what we guess)
\end{itemize}

\noindent
When truth values are primary, these properties are difficult or
impossible to achieve.  Different truth-value schemes
(strength/confidence, fuzzy logic, interval arithmetic, probability
values) require different combination rules, different query
algorithms, and different validation approaches.  The result is
fragmentation.  The world-model calculus reunifies this fragmented
landscape by locating the common structure beneath the diverse surface
phenomena.

What remains is engineering, extension, and empirical validation.
The mathematical foundation is in place.

\section*{In Summary}

The world-model calculus is a calculus of revisable uncertain state.
A system stores what it knows in a posterior state, not in isolated
truth values.  Queries ask questions of that state, and the raw
answers are evidence values.  Strength, confidence, and intervals are
useful summaries---views extracted for display, thresholding, and
action---but they are not themselves the thing being stored or proved.

The world-model calculus does not replace existing reasoning systems.
It provides the semantic foundation they can share.  A ProbLog
program, a Bayesian network, and a PLN rule base are no longer
separate formalisms requiring separate toolchains.  They are different
languages for specifying world models, and they share common semantics
for revision, extraction, and view projection.

This is the unity that reasoning systems have long needed: not the
false unity of a single approach that claims to solve everything, but
the true unity of a common foundation that reveals why different
approaches work and how they relate.

\begin{quote}
\emph{Probability is what you query.  Evidence is what you carry.
Everything else follows.}
\end{quote}
